\section{Related Works}
\label{sec:related}
\if0
NetDream sits at the intersection of Kubernetes autoscaling and learned
control.  The closest prior systems recognize either that autoscaling
should be predictive, or that microservice systems should be modeled as
graphs.  Our distinction is that NetDream uses a graph-structured world
model not only to predict a future metric, but to plan over imagined
future system trajectories and reject unsafe actions before execution.

\subsection{From Reactive Autoscaling to Predictive Kubernetes Control}
\label{sec:related-reactive-predictive}

Production Kubernetes autoscaling is still dominated by reactive control.
The Horizontal Pod Autoscaler (HPA) scales replica counts when CPU
utilization or custom metrics cross a configured threshold
\citep{rzadca2020autopilot,googlecloud2025hpa,aws2025hpa}.  This design
is simple, robust, and easy to operate, which explains its broad
adoption.  Its limitation is also fundamental: HPA observes pressure only
after it appears.  It cannot predict whether scaling one service will
shift load to another service, reduce downstream latency, or create a
future service-level violation.

A line of work augments HPA with forecasting or adaptive rules.
SLA-adaptive HPA dynamically adjusts thresholds
\citep{pozdniakova2024sla}; Predictive HPA forecasts demand several
minutes ahead \citep{mondal2023predictive}; Smart HPA uses recurrent
time-series models \citep{ahmad2024smarthpa}; TAHPA exploits diurnal
patterns \citep{phuc2022tahpa}; and Gwydion switches between reactive
and predictive modes \citep{santos2024gwydion}.  These methods improve
when the workload has predictable temporal structure, but they remain
limited by the same control interface: forecast a metric, then trigger a
rule.  They do not learn a transition model of the microservice graph,
nor do they simulate the multi-step consequences of candidate scaling
actions before acting.

NetDream takes a different view.  Autoscaling is not merely forecasting
future load; it is choosing among counterfactual futures.  The controller
must ask what will happen if it scales the frontend, the checkout
service, or the recommendation service, and how each action will affect
latency, cost, and violations over the next several control steps.
NetDream addresses this by learning an action-conditioned graph world
model and using it inside a planner.  This turns predictive autoscaling
from metric forecasting into model-based decision making.

\subsection{Learning-Based Autoscaling: From Metric Prediction to World Models}
\label{sec:related-learning-autoscaling}

\paragraph{Reactive Kubernetes autoscaling.}
Production autoscaling is still dominated by Kubernetes' Horizontal Pod
Autoscaler (HPA), which scales replicas once CPU or custom-metric
utilization crosses a configured threshold
\citep{rzadca2020autopilot,googlecloud2025hpa,aws2025hpa}.
HPA is simple and operationally proven but cannot anticipate downstream
load, predict latency cascades, or simulate the multi-step consequences
of a scaling decision.
Variants that augment HPA with short-horizon demand forecasting
\citep{mondal2023predictive,ahmad2024smarthpa,santos2024gwydion} or
adaptive thresholds
\citep{pozdniakova2024sla,phuc2022tahpa}
improve responsiveness when traffic has predictable temporal structure,
but they still trigger scaling rules from a forecast metric rather than
from a simulated future trajectory.

\paragraph{Topology-aware prediction and model-free control.}
\fi

Learning-based autoscaling has moved beyond hand-designed rules, but
most methods still fall into two categories: predict a metric and react,
or learn a policy directly.  GNN-based autoscalers exploit the fact that
microservice systems are graphs.  GRAF predicts tail latency with a GNN
and uses the prediction for proactive SLO-aware allocation
\citep{park2021graf}.  pHPA continuously trains a GNN to infer tail
latency along microservice chains \citep{choi2021phpa}.  DeepScaler
combines spatio-temporal GNNs with adaptive graph learning for holistic
autoscaling \citep{meng2023deepscaler}; Graph-PHPA combines LSTM and GNN
modules for proactive horizontal scaling \citep{dang2022graphphpa}; AGQ
uses dense spatio-temporal GNNs to capture indirect dependencies
\citep{liang2025agq}; and USRFNet learns unified system representations
for microservice tail-latency prediction \citep{zhang2025usrfnet}.  These
systems show that topology-aware representations matter, but their
control logic remains prediction-driven rather than imagination-driven:
the GNN predicts a future metric, and a separate heuristic or policy
decides how to scale. Reinforcement-learning autoscalers instead learn control policies
directly.  Prior work has explored DQN, PPO, TD3, and recurrent
reinforcement learning for microservice and serverless autoscaling
\citep{abdelkhaleq2023dqn,qiu2023aware,zhang2025kiss,drpc2024,
agarwal2024deep}.  Multi-objective variants optimize SLO and cost
trade-offs through bi-level control, Lagrangian relaxation, or Pareto
optimization \citep{wang2024autothrottle,wang2024deepscaling,
marchese2025slocost}.  GraphPilot is closest in spirit among GNN-based
controllers: it combines a Temporal Graph Network encoder with an
actor--critic policy for Kubernetes autoscaling
\citep{chen2025graphpilot}.  However, these approaches learn a policy,
not a graph world model used for explicit planning.  They do not roll out
candidate action sequences in a learned dynamics model and reject unsafe
futures before execution. Digital-twin approaches move closer to model-based control. KARMA builds
a neural transition model of the cluster and trains multi-agent policies
in simulation for resilience against failures such as DDoS attacks and
pod crashes \citep{karma2025}.  However, its transition model is not
graph-structured, and its objective is adversarial resilience rather than
cost--SLO autoscaling on standard microservice workloads. 
NetDream differs in three ways: it learns topology-aware graph dynamics, uses the learned model online for model-predictive planning, and applies an
explicit safety filter that rejects candidate action sequences predicted
to violate service-level constraints. 
In summary, prior Kubernetes autoscalers either react to observed
pressure, forecast a metric and apply a rule, or learn a model-free
policy. NetDream instead treats autoscaling as planning with a learned
graph world model.  This distinction is the source of its empirical
advantage: the controller can compare multiple possible futures, reason
over service dependencies, and avoid unsafe plans before they are applied
to the real cluster.